% Einheit 1 von 3 – Modell und Glockenkurve: die Dichtefunktion lesen und skizzieren (μ als Symmetrieachse und Maximumsstelle, σ als Breitenmaß (bank/normalverteilung-und-sigma-regeln/e1.jsonl, zusammenbau v0.1)
%% TODO Zeitmarke Einheit 1: Marken-Zeile nicht lesbar
%% TODO Zweigzeile Teil 1: was hier gelernt wird (Satz in Schülersprache aus dem Einheitstitel) – Einheit 1
\einheitenkopf[e1]{Einheit 1 von 3 · Modell und Glockenkurve: die Dichtefunktion lesen und skizzieren (μ als Symmetrieachse und Maximumsstelle, σ als Breitenmaß}
\zweigzeile{keine Prüfungsaufgabe}
\setcounter{aufgabe}{7}

%% TODO Titel: Ich-kann-Satz fehlt in der Bank (Platzhalter: Kettenname) – Nr. 8
%% TODO Anweisung: ein Satz über der ersten Teilaufgabe fehlt in der Bank – Nr. 8
\begin{aufgabe}{Modell und Glockenkurve}
%% TODO \rechenplatz{n}: Schrittzahl des Grundfalls fehlt in der Bank (nur bei fester Schreibform, 2.3 b)
\begin{teile}
\teil Wie wahrscheinlich wiegt ein Brot zwischen $480$ g und $520$ g? Gemeint ist … \\ \kreuz{Fläche unter der Dichte} \\ \kreuz{Höhe des Graphen}
\teil Die Abbildung zeigt die Dichte von $X$ (die Masse eines Brötchens in g) mit $\mu = 50$ und $\sigma = 5$. Beschreibe die Wirkung von $\sigma = 10$ bei gleichem $\mu$ und skizziere die neue Dichte. \hfill (Abitur 2022 LK)

\normalverteilung{50}{5}
\teil Die Abbildung zeigt die Dichte einer normalverteilten Zufallsgröße $X$ mit dem Erwartungswert $40$. Wie groß ist $P(X = 35)$? \hfill (Abitur 2024 LK) \\

\normalverteilung{40}{4}
\leerfeld
\teil Die Füllmenge $X$ einer Flasche in ml ist normalverteilt mit $\mu = 500$ und $\sigma = 4$. Im Modell ist $P(X < 0) > 0$. Ist das ein Argument gegen das Modell? \hfill (Abitur 2026 LK)
\teil Eine Anlage füllt Mehlpackungen; die Füllmenge $X$ in g ist normalverteilt mit $\mu = 1003$ und $\sigma = 4$ (Abbildung). Vorschlag 1: die Füllmenge im Mittel erhöhen. Vorschlag 2: genauer abfüllen. Skizziere zu jedem Vorschlag eine passende Dichte und begründe, dass $P(X < 1000)$ jeweils kleiner wird. \hfill (Abitur 2023 LK)

\normalverteilung[bis=1000]{1003}{4}
\end{teile}
\end{aufgabe}

%% TODO Titel: Ich-kann-Satz fehlt in der Bank (Platzhalter: Kettenname) – Nr. 9
%% TODO Anweisung: ein Satz über der ersten Teilaufgabe fehlt in der Bank – Nr. 9
\begin{aufgabe}{Modell und Glockenkurve – Fehler finden, Begründen}
\begin{teile}
\teil Ich finde den Fehler: Die Dichte von $X$ hat an der Stelle $20$ den Wert $0{,}08$. Tim: „Also ist $P(X = 20) = 0{,}08$.“
\teil Begründe: Die Fläche unter jeder Dichtefunktion ist eins.
\end{teile}
\end{aufgabe}
